% The Analytic Addition Formula  (source §18; added: an independent proof of the
% wp addition theorem from Abel + the ODE, the determinant form, the reflected
% sign via oddness of wp', and the duplication formula)
\chapter{The Analytic Addition Formula}
\label{ch:analytic_addition}

The chord computation of \cref{ch:coordinate_formulas} has an analytic
counterpart: a closed formula for \(\wp(u+v)\) in terms of \(\wp\) and \(\wp'\) at
\(u\) and \(v\). This chapter proves the Weierstrass addition theorem directly
from Abel's law and the differential equation --- independently of the geometry ---
and then shows it reproduces the chord law exactly, with the oddness of \(\wp'\)
accounting for the reflection. The two derivations meeting is the analytic seal on
\cref{thm:main}.

\section{The Addition Theorem}
\label{sec:aa-theorem}

\begin{theorem}[Weierstrass addition formula]
\label{thm:aa-addition}
For \(u, v \notin \Lambda\) with \(u \not\equiv \pm v \pmod\Lambda\),
\[
  \boxed{\,\wp(u+v) = -\wp(u) - \wp(v)
  + \frac14\left( \frac{\wp'(u) - \wp'(v)}{\wp(u) - \wp(v)} \right)^2 .\,}
\]
\end{theorem}

\begin{proof}
Let \(A = \dfrac{\wp'(u) - \wp'(v)}{\wp(u) - \wp(v)}\) and \(B = \wp'(u) - A\wp(u)\),
so that the line \(Y = A\,X + B\) passes through \((\wp(u), \wp'(u))\) and
\((\wp(v), \wp'(v))\). Consider the elliptic function
\[
  h(z) = \wp'(z) - A\,\wp(z) - B .
\]
It is elliptic with a single pole at \(0\), of order three (from \(\wp'\)), so it
has order three and, by \cref{thm:ell-valence}, exactly three zeros counted with
multiplicity. Two of them are \(u\) and \(v\) (where the line meets the curve);
let \(w\) be the third. By Abel's law (\cref{thm:ell-abel}) applied to
\(\divisor(h) = [u] + [v] + [w] - 3[0]\),
\[
  u + v + w \equiv 0 \pmod\Lambda, \qquad\text{so}\qquad w \equiv -(u+v).
\]

At a zero \(z\) of \(h\) we have \(\wp'(z) = A\wp(z) + B\); squaring and using the
differential equation \(\wp'^2 = 4\wp^3 - g_2\wp - g_3\) (\cref{thm:wp-ode}),
\[
  4\,\wp(z)^3 - g_2\,\wp(z) - g_3 = \bigl(A\,\wp(z) + B\bigr)^2
  = A^2 \wp(z)^2 + 2AB\,\wp(z) + B^2 .
\]
Thus \(\wp(z)\) is a root of the cubic \(4X^3 - A^2 X^2 - \cdots = 0\), whose three
roots are \(\wp(u), \wp(v), \wp(w)\). Vieta's formula for the sum of the roots
gives
\[
  \wp(u) + \wp(v) + \wp(w) = \frac{A^2}{4} .
\]
Finally \(\wp(w) = \wp(-(u+v)) = \wp(u+v)\) because \(\wp\) is even
(\cref{prop:wp-even-periodic}). Solving,
\(\wp(u+v) = \tfrac{A^2}{4} - \wp(u) - \wp(v)\), which is the boxed formula.
\end{proof}

The proof used only Abel's law, the differential equation, and the evenness of
\(\wp\) --- no reference to the cubic in \(\CP^2\). It nevertheless reconstructs the
chord rule.

\section{Agreement with the Chord Law}
\label{sec:aa-chord}

Set \(P = \Phi(u) = (\wp(u), \wp'(u))\), \(Q = \Phi(v) = (\wp(v), \wp'(v))\), and
\(x_1 = \wp(u)\), \(y_1 = \wp'(u)\), \(x_2 = \wp(v)\), \(y_2 = \wp'(v)\). The chord
slope is
\[
  m = \frac{y_2 - y_1}{x_2 - x_1}
    = \frac{\wp'(v) - \wp'(u)}{\wp(v) - \wp(u)}
    = \frac{\wp'(u) - \wp'(v)}{\wp(u) - \wp(v)} = A .
\]
The chord formula of \cref{thm:coord-add} gives \(x_{P\oplus Q} = \tfrac{m^2}{4} -
x_1 - x_2\), and the addition theorem gives \(\wp(u+v) = \tfrac{A^2}{4} - \wp(u) -
\wp(v)\); since \(m = A\), these are identical:
\[
  x_{P \oplus Q} = \wp(u+v).
\]
The two computations --- summing the roots of the intersection cubic, and summing
the values of \(\wp\) at the zeros of \(h\) --- are the same calculation viewed
through \(\Phi\).

\section{Why the \texorpdfstring{$y$}{y}-Coordinate Has the Reflected Sign}
\label{sec:aa-reflection}

The third intersection of the chord is \(R = \Phi(w) = \Phi(-(u+v))\). Its
coordinates are
\[
  R = \bigl( \wp(-(u+v)),\, \wp'(-(u+v)) \bigr)
    = \bigl( \wp(u+v),\, -\wp'(u+v) \bigr),
\]
using \(\wp\) even and \(\wp'\) odd. Reflecting,
\[
  \ominus R = \bigl( \wp(u+v),\, \wp'(u+v) \bigr) = \Phi(u+v) = P \oplus Q .
\]
So the \(y\)-coordinate of the sum is \(+\wp'(u+v)\), while the third intersection
had \(-\wp'(u+v)\): the reflection across the \(x\)-axis is exactly the oddness of
\(\wp'\).

\begin{keyidea}{Analytic and geometric agreement}
The addition formula for \(\wp\) is the analytic form of the chord law. The
root-sum computation produces the same \(x\)-coordinate as \(\wp(u+v)\), and the
oddness of \(\wp'\) supplies the reflection that produces the \(y\)-coordinate.
The minus sign in ``reflect \(R\)'' and the minus sign in \(\wp'(-z) = -\wp'(z)\)
are one and the same.
\end{keyidea}

\section{The Determinant Form}
\label{sec:aa-determinant}

Collinearity of \(P\), \(Q\), and \(R\) is a vanishing determinant, which packages
the addition theorem symmetrically.

\begin{proposition}
\label{prop:aa-determinant}
For \(u, v\) generic,
\[
  \det\begin{pmatrix}
    1 & \wp(u)   & \wp'(u) \\
    1 & \wp(v)   & \wp'(v) \\
    1 & \wp(u+v) & -\wp'(u+v)
  \end{pmatrix} = 0 .
\]
\end{proposition}

\begin{proof}
The three rows are the coordinates \((1, x, y)\) of \(P = \Phi(u)\),
\(Q = \Phi(v)\), and \(R = \Phi(-(u+v)) = (\wp(u+v), -\wp'(u+v))\). By
\cref{thm:lin-collinear} these three points are collinear --- they are the
intersection of the chord with \(E\) --- and three points of the affine plane are
collinear precisely when this determinant vanishes.
\end{proof}

Expanding the determinant and solving for \(\wp(u+v)\) returns
\cref{thm:aa-addition}; the symmetric form makes the role of all three points
\(u, v, -(u+v)\) manifest.

\section{The Duplication Formula}
\label{sec:aa-duplication}

Letting \(v \to u\) gives the formula for \(\wp(2u)\).

\begin{corollary}[Duplication]
\label{cor:aa-duplication}
For \(u\) with \(\wp'(u) \ne 0\),
\[
  \wp(2u) = -2\wp(u) + \frac14\left( \frac{\wp''(u)}{\wp'(u)} \right)^2
          = -2\wp(u) + \frac14\left( \frac{6\wp(u)^2 - \tfrac12 g_2}{\wp'(u)} \right)^2 .
\]
\end{corollary}

\begin{proof}
As \(v \to u\), the slope \(A = \dfrac{\wp'(u) - \wp'(v)}{\wp(u) - \wp(v)}\) is a
difference quotient; by l'H\^opital (dividing numerator and denominator by
\(u - v\) and taking the limit) it tends to \(\dfrac{\wp''(u)}{\wp'(u)}\).
\Cref{thm:aa-addition} then gives \(\wp(2u) = -2\wp(u) + \tfrac14(\wp''/\wp')^2\).
Differentiating \(\wp'^2 = 4\wp^3 - g_2\wp - g_3\) yields \(\wp'' = 6\wp^2 -
\tfrac12 g_2\) (\cref{ex:wp-recurrence}), giving the second form.
\end{proof}

This matches the doubling formula of \cref{thm:coord-double}: the tangent slope
\((12 x^2 - g_2)/(2y)\) is exactly \((6\wp^2 - \tfrac12 g_2)/\wp' =
\wp''/\wp'\), with \(x = \wp(u)\), \(y = \wp'(u)\).

\section{Summary}
\label{sec:aa-summary}

The Weierstrass addition formula (\cref{thm:aa-addition}) gives \(\wp(u+v)\) as
\(-\wp(u) - \wp(v)\) plus a quarter-square of the slope \(A\); it is proved
directly from Abel's law and the differential equation, by summing the values of
\(\wp\) at the three zeros of \(h = \wp' - A\wp - B\). It reproduces the chord
\(x\)-coordinate, and the oddness of \(\wp'\) supplies the reflected \(y\)-sign
(\cref{sec:aa-chord,sec:aa-reflection}). In symmetric form it is a vanishing
\(3 \times 3\) determinant (\cref{prop:aa-determinant}), and the \(v \to u\) limit
is the duplication formula (\cref{cor:aa-duplication}), matching the geometric
doubling. Analytic and geometric addition are one operation.

\section*{Exercises}
\addcontentsline{toc}{section}{Exercises}

\begin{exercise}
\label{ex:aa-h-order}
Verify the claims about \(h(z) = \wp'(z) - A\wp(z) - B\) in the proof of
\cref{thm:aa-addition}: that it is elliptic of order three with a triple pole at
\(0\), and that \(u, v\) are among its zeros. Why does the case
\(u \equiv \pm v\) have to be excluded?
\end{exercise}

\begin{exercise}
\label{ex:aa-determinant-expand}
Expand the determinant of \cref{prop:aa-determinant} along the first column and
solve the resulting linear relation for \(\wp(u+v)\), recovering
\cref{thm:aa-addition}.
\end{exercise}

\begin{exercise}
\label{ex:aa-duplication-match}
Show that the duplication slope \(\wp''(u)/\wp'(u)\) equals the tangent slope
\((12x^2 - g_2)/(2y)\) of \cref{thm:coord-double} with \(x = \wp(u)\),
\(y = \wp'(u)\), confirming \cref{cor:aa-duplication} matches the geometric
doubling formula.
\end{exercise}

\begin{exercise}
\label{ex:aa-symmetry}
Check that \cref{thm:aa-addition} is symmetric under \(u \leftrightarrow v\), and
compute the limit \(v \to 0\) to confirm consistency with \(\wp(u+0) = \wp(u)\).
(The pole of \(\wp(v)\) at \(v = 0\) cancels against the slope term.)
\end{exercise}

\begin{exercise}
\label{ex:aa-three-term}
Apply \cref{thm:aa-addition} to \(\wp(u-v)\) by replacing \(v\) with \(-v\) (so the
slope becomes \((\wp'(u) + \wp'(v))/(\wp(u) - \wp(v))\), using \(\wp\) even,
\(\wp'\) odd). Adding the formulas for \(\wp(u+v)\) and \(\wp(u-v)\), derive
\[
  \wp(u+v) + \wp(u-v) = -2\wp(u) - 2\wp(v)
  + \frac{\wp'(u)^2 + \wp'(v)^2}{2\,(\wp(u) - \wp(v))^2},
\]
and conclude, using the differential equation, that the right-hand side is a
rational function of \(\wp(u)\) and \(\wp(v)\).
\end{exercise}
